% !TEX TS-program = lualatex
Необходимо разработать программный продукт для игры в нарды на персональном компьютере через веб-браузер. Игра должна поддерживать два режима:

\begin{enumerate}
    \item Человек против компьютера.
    \item Компьютер против компьютера (интеллектуального агента).
\end{enumerate}

Для каждого из режимов требуется разработать алгоритм интеллектуального агента, который позволяет гибко управлять уровнем сложности игры. Предполагается, что агент будет создан на основе модели машинного обучения. Требования к интеллектуальному агенту:

\begin{enumerate}
    \item Возможность обучения интеллектуального агента на персональном компьютере пользователя.
    \item Возможность настройки уровня сложности искусственного интеллекта с помощью определённых параметров.
\end{enumerate}

Для режима <<человек против компьютера>> необходимо реализовать удобный пользовательский интерфейс (\refimage{backgammon-ui}), а также специальное программное обеспечение для:

\begin{enumerate}
    \item Управления пользовательскими партиями.
    \item Загрузки противников (компьютерных агентов) в игру.
\end{enumerate}

Для режима <<компьютер против компьютера>> требуется разработать специальное программное обеспечение, предназначенное для тренировок компьютерных агентов путём автоматической игры друг с другом.

Требования к программному обеспечению:

\begin{enumerate}
    \item Программное обеспечение должно включать следующие компоненты:
    \begin{itemize}
        \item серверную часть;
        \item клиентскую часть;
        \item инструменты для работы с интеллектуальными агентами (тренировка, загрузка).
    \end{itemize}
    \item Приложение должно функционировать в сети интернет.
    \item Приложение не должно сохранять данные о пользователях, при этом оно должно быть устойчиво к попыткам нечестной игры.
\end{enumerate}

\image[width=0.9\textwidth]{Интерфейс клиентской части приложения}{backgammon-ui}{backgammon-ui}
